\begin{figure}[htbp]
\centering
\textbf{Frequency Spectrum of a Damped Sinusoid}

\vspace{0.3cm}

% Time domain signal
\begin{tikzpicture}
\begin{axis}[
    width=13cm, height=4.5cm,
    xlabel={Time $t$ (seconds)},
    ylabel={$f(t)$},
    xmin=0, xmax=5,
    ymin=-1.1, ymax=1.1,
    grid=major,
    grid style={gray!30},
    title={\normalsize Time Domain: $f(t) = \ee^{-\sigma t}\cos(\omega_0 t)$ \quad ($\sigma = 0.5$, $\omega_0 = 4\pi$ rad/s)},
    title style={at={(0.5,1.05)}, anchor=south}
]

% Damped sinusoid: exp(-0.5*t)*cos(4*pi*t) = exp(-0.5*t)*cos(2*pi*2*t), f_0 = 2 Hz
\addplot[UofSCAtlantic, very thick, domain=0:5, samples=500]
    {exp(-0.5*x)*cos(deg(4*pi*x))};

% Envelope (positive)
\addplot[UofSCGarnet, dashed, thick, domain=0:5, samples=100]
    {exp(-0.5*x)};

% Envelope (negative)
\addplot[UofSCGarnet, dashed, thick, domain=0:5, samples=100]
    {-exp(-0.5*x)};

% Annotation
\node[UofSCGarnet, font=\small] at (3.2, 0.35) {Envelope: $\ee^{-\sigma t}$};
\draw[UofSCGarnet, thin, ->] (2.7, 0.3) -- (2.0, 0.38);

\node[UofSCAtlantic, font=\small] at (1.5, -0.65) {$f_0 = 2$ Hz};

\end{axis}
\end{tikzpicture}

\vspace{0.6cm}

% Magnitude spectrum
\begin{tikzpicture}
\begin{axis}[
    width=6.2cm, height=5cm,
    xlabel={Frequency $f$ (Hz)},
    ylabel={$|F(\jj\omega)|$},
    xmin=-6, xmax=6,
    ymin=0, ymax=1.2,
    grid=major,
    grid style={gray!30},
    title={\normalsize Magnitude Spectrum},
    title style={at={(0.5,1.05)}, anchor=south},
    xtick={-6,-4,-2,0,2,4,6},
    ytick={0, 0.5, 1}
]

% Magnitude spectrum: |F(jw)| = 0.5 * [sigma/((sigma)^2 + (w-w0)^2) + sigma/((sigma)^2 + (w+w0)^2)]
% For damped cosine: peaks at +/- f_0 with width determined by sigma
% Using Lorentzian-like shape centered at f0=2 Hz

% Peak at +f_0 = 2 Hz
\addplot[UofSCAtlantic, very thick, domain=-6:6, samples=300]
    {0.5*0.5/((0.5)^2 + (2*pi*(x-2))^2) + 0.5*0.5/((0.5)^2 + (2*pi*(x+2))^2)};

% Mark the center frequencies
\draw[UofSCGarnet, dashed, thick] (2, 0) -- (2, 1.1);
\draw[UofSCGarnet, dashed, thick] (-2, 0) -- (-2, 1.1);

% Annotations
\node[UofSCGarnet, font=\footnotesize, anchor=south] at (2, 1.12) {$f_0$};
\node[UofSCGarnet, font=\footnotesize, anchor=south] at (-2, 1.12) {$-f_0$};

% Half-power bandwidth annotation
\draw[UofSCHorseshoe, thick, <->] (1.84, 0.5) -- (2.16, 0.5);
\node[UofSCHorseshoe, font=\tiny, anchor=west] at (2.2, 0.5) {$\Delta f$};

\end{axis}
\end{tikzpicture}
\hfill
% Phase spectrum
\begin{tikzpicture}
\begin{axis}[
    width=6.2cm, height=5cm,
    xlabel={Frequency $f$ (Hz)},
    ylabel={$\angle F(\jj\omega)$ (degrees)},
    xmin=-6, xmax=6,
    ymin=-100, ymax=100,
    grid=major,
    grid style={gray!30},
    title={\normalsize Phase Spectrum},
    title style={at={(0.5,1.05)}, anchor=south},
    xtick={-6,-4,-2,0,2,4,6},
    ytick={-90, -45, 0, 45, 90},
    yticklabels={$-90$, $-45$, $0$, $45$, $90$}
]

% Phase spectrum: transitions through the resonance frequencies
% Phase = atan2(Im, Re) for the transform
% Near +f_0: phase transitions from +90 to -90
% Near -f_0: phase transitions from -90 to +90

% Approximate phase behavior (simplified for visualization)
\addplot[UofSCCongaree, very thick, domain=-6:-2.3, samples=50] {-90 + 90*2/pi*atan(deg(5*(x+2)))};
\addplot[UofSCCongaree, very thick, domain=-1.7:1.7, samples=50] {0};
\addplot[UofSCCongaree, very thick, domain=1.7:6, samples=50] {90 - 90*2/pi*atan(deg(5*(x-2)))};

% Transition regions
\addplot[UofSCCongaree, very thick, domain=-2.3:-1.7, samples=30]
    {-90*tanh(3*(x+2))};
\addplot[UofSCCongaree, very thick, domain=1.7:2.3, samples=30]
    {-90*tanh(3*(x-2))};

% Mark center frequencies
\draw[UofSCGarnet, dashed, thick] (2, -100) -- (2, 100);
\draw[UofSCGarnet, dashed, thick] (-2, -100) -- (-2, 100);

% Annotations
\node[UofSCGarnet, font=\footnotesize, anchor=south] at (2, 92) {$f_0$};
\node[UofSCGarnet, font=\footnotesize, anchor=south] at (-2, 92) {$-f_0$};

\end{axis}
\end{tikzpicture}

\vspace{0.4cm}

% Discrete spectrum example (for periodic signal)
\begin{tikzpicture}
\begin{axis}[
    width=13cm, height=4cm,
    xlabel={Harmonic Number $n$},
    ylabel={$|c_n|$},
    xmin=-0.5, xmax=12.5,
    ymin=0, ymax=1.4,
    grid=major,
    grid style={gray!30},
    title={\normalsize Discrete Magnitude Spectrum of Square Wave (Fourier Series Coefficients)},
    title style={at={(0.5,1.05)}, anchor=south},
    xtick={0,1,2,3,4,5,6,7,8,9,10,11,12},
    ytick={0, 0.5, 1, 1.27},
    yticklabels={$0$, $0.5$, $1$, $4/\pi$}
]

% Magnitude of Fourier coefficients for square wave: |c_n| = 4/(n*pi) for odd n, 0 for even n
% Plot as stem plot using ycomb

\addplot+[UofSCAtlantic, very thick, ycomb, mark=*, mark options={fill=UofSCAtlantic, scale=1.2}]
    coordinates {
        (0, 0)
        (1, 1.273)
        (2, 0)
        (3, 0.424)
        (4, 0)
        (5, 0.255)
        (6, 0)
        (7, 0.182)
        (8, 0)
        (9, 0.141)
        (10, 0)
        (11, 0.116)
        (12, 0)
    };

% Envelope curve (1/n decay)
\addplot[UofSCGarnet, dashed, thick, domain=0.8:12, samples=50]
    {4/(pi*x)};

% Annotations
\node[UofSCGarnet, font=\footnotesize] at (8, 0.65) {Envelope: $4/(n\pi)$};
\node[UofSCAtlantic, font=\footnotesize] at (6.5, 0.12) {Even harmonics = 0};

\end{axis}
\end{tikzpicture}

\caption{Frequency domain representation of signals. \textbf{Top:} Time-domain damped sinusoid $f(t) = \ee^{-\sigma t}\cos(\omega_0 t)$ with decay rate $\sigma = 0.5$ s$^{-1}$ and frequency $f_0 = 2$ Hz. The dashed Garnet lines show the exponential envelope. \textbf{Middle:} Magnitude spectrum (Atlantic blue) shows peaks at $\pm f_0$ with bandwidth $\Delta f$ proportional to $\sigma$. Phase spectrum (Congaree teal) shows $\pm 90°$ phase transitions at the resonance frequencies. \textbf{Bottom:} Discrete spectrum for a periodic square wave showing only odd harmonics with magnitude decreasing as $4/(n\pi)$.}
\label{fig:spectrum_visualization}
\end{figure}
